\documentclass[12pt]{article}

\usepackage[a4paper,margin=1in]{geometry}
\usepackage{graphicx}
\usepackage{setspace}

\begin{document}

\begin{titlepage}
    \centering

    % School Logo
    \includegraphics[width=0.40\textwidth]{logo.png}\par\vspace{1cm}

    % School Name
    {\Large Gisma University of Applied Science\par}
    \vspace{2cm}

    % Report Title
    {\Huge\bfseries A Software Engineer in the Making \par}
    \vspace{2cm}

    % Student Information
    {\large
    \textbf{Prepared by:}\\
    Princia Mirielle Ishimwe\\[0.5cm]
    
    \textbf{Student ID:}\\
    GH1049008
    }

    \vfill

    % Course Information (Optional)
    {\large
    Computer Science lab\\
    Teacher: William Baker Morrison
    
    
    }

    \vspace{1cm}

    % Date
    {\large \today\par}

\end{titlepage}

\newpage
\section{Declaration}
I, Princia Mirielle Ishimwe, declare that this report titles "A software engineer in the making" is my own original work and has not been submitted elsewhere in my academic journey.





Signed: Princia Mirielle Ishimwe 

Date: 24.06.2026

\newpage



\tableofcontents
\newpage


%---------------- Table of Contents ----------------%

\section{Abstract and Acknowledgement}
\subsection{Abstract}

This is a report about the development of an online portfolio of Princia Mirielle Ishimwe. Different structures were used to develop the portfolio, the CV in it and this report. The design of each and every document was chosen carefully and with a genuine thought behind. Tools and technologies like Overleaf, Git, Github, Latex were used to develop this report and everything needed to complete it. Strength and weaknesses of the portfolio were shown in this report as well as future recommendations. The report reveals the skills and different levels of Princia Mirielle Ishimwe, as an aspiring software engineer.

\paragraph{Key words: Portfolio, Website, Github, repositories, Latex, Overleaf, HTML/CSS}


\subsection{Acknowledgement}

I would like to wholeheartdly thank our module tutor, Mr William Baker Morrison, for his valuable insights, his guidance and his teaching skills shown throughout the development of this report, secondly I thank Gisma University of Applied Sciences for the modules they prepare for us, as well as the learning methods to acquire professional and soft skills at the same time , and lastly I thank with all my heart my family and my friends for their constant support, prayers and motivation that keep me grounded and encouraged.

\newpage

%---------------- Introduction ----------------%
\bfseries\section {Introduction}

Before the centralization and digitalisation of resources, job seekers were relying mainly on their resumes and especially hand-written letters from previous employers, church leader or community leaders( Insler, 2026). In the early 1900s, following the launch of the first online job board, job seekers started to convert their works into digital formats (Bennett, nd). In the late 2000s because of the expandure of the internet and launch of Linkedin, online portfolio were more introducted and used.(Cait, 2023).

\paragraph{An electric or online portfolio is a digital collection of informations that showcases someone's skills, experiences and achievements over time. (Cait, 2023)}

\paragraph{This report is about the development of an online portfolio of me, Princia Mirielle Ishimwe, it starts with structure and links of my portfolio and Github, secondly, it explains the design decisions I made while developing the portfolio and also the Latex Cv in it, thirdly, it explains the tools and technologies used to develop my portfolio, my CV as well as this report, and lastly the elements I found as strengths and weeknesses of my portfolio are explained without forgetting the recommendations for the development of my future portfolio.  }

\paragraph{ The first chapter is about the links to my online portfolio and Github repositories containing my CV, some technical exercises and this report}




\newpage

%---------------- Links ----------------%
\bfseries\section{Online Portfolio and Github repository links}

Here are the links:


For my online portfolio: \href{ https://mirielle-dev.github.io/Mirielle.github.io/}

For my LateX CV:  \href{https://github.com/Mirielle-dev/latex-cv.git}

For my technical exercise: \href{https://github.com/Mirielle-dev/My-python-projects } 
(see appendix 2)

For my report: 








%---------------- Structure ----------------%
\bfseries\section{The Structure of my Portfolio and Github Repository}
\bfseries\subsection{The Structure of my Online Portfolio}
 The structure of my online portfolio was intentionally guided towards the objective to introduce who I am, my strengths,my projects and my goals. Certainly, it was done according to some examples shown in class and research I've made about the look and content of a professional portfolio.

 
 \paragraph{According to Larousse Dictionary, an online portfolio is a digital showcase of someone's work, skills and experiences.
 According to that I choose to divide my online portfolio in 4 main parts:
 \textit{Introduction page, About me page, Skills page, Projects page and Contacts page}. }

 \bfseries\subsubsection{Introduction page}

 \textit{An introduction page} or \textit{Home page} is the first page setting the tone of the following pages, it is made to attract the reader's attention, as Botany Suits said \textit{"You never get a second chance to make a first impression"}, so the first impression matter.


\paragraph{For my \textit{Home page}, I highlighted it with a picture of mine not so casual and not so professional and with a big smile, and I added shortcut buttons to get access to my resume and every other page of my portfolio. My full name was clearly highlighted and who I am and my daily motivational quote.
I truly believe it represents who I am and someone reading can get an idea of what I'm intending to.}


\subsubsection{About me page} 
\textit{About me} section is the section that introduces a person, shows their background and their personality. For me, the \textit{About me} section highlighted the skills I'm building and the ones I'm intending to learn, as a consistent learner, showing who I am and how I am as a person and the way I'm adjusting to this beautiful country while learning the language. 

\subsubsection{Skills page}
\textit{Skills page} is a dedicated page in the portfolio to offshow your capacities, your expertise and skills professional ones as well as soft ones.(The Learning Portal,2022)
For my case as an aspiring software engineer, I choose to list the skills from what I'm learning right now to what I learned throughout my classes and even in High school and highlighted them using arranged cases.

\subsubsection{Project pages}
\subsubsection{Contact page}
\textit{The contact page} is the page with your contacts, as for the aspiring software engineer I am, I choose to put there my email for personal or preference query, my linkedin to connect with potential employers and my Github to offshow my projects and skills applied.

\subsection{The Structure of my GitHub Repository}

The structure of my Github repositories is based mainly on what I needed to deliver for this assignment, I have 4 main repositories and the benefits of separating content into folders is to make updates simpler and provide more readability.

\paragraph {* Mirielle.github.io repository: This is the repository for my online portfolio. It contains my HTML/CSS code for my portfolio, all the picture and Cv used for the completion of my portfolio.}

\paragraph {* Latex-CV repository: This repository contains my Latex-CV with both a .pdf and .tex file to ensures transparency and compatibility between my CV and my online portfolio.}

\paragraph {* My Python Project repository: As an aspiring software engineer, I chose to start my programming journey with Python, even if I did C++ and SQL back in HS, this repository contains my python projects and shows my willngness to learn even if the projects are small.}

\paragraph {* Report repository: This repository contains my report "A software engineer in the making" and it includes a .pdf and .tex file. }

\newpage


%---------------- Design decisions ----------------%
\section{Design Decisions}

\paragraph { \textit{Blue} means and offshows loyalty, care, responsability, reliability, peace and calm, and \textit{Black} means authority, power, mystery and serious in business and technology(Cerrato, 2012). For me, choosing Blue and Black as main colors of my portfolio was based on these facts and Blue is from my favorite cartoon Vaiana (See appendix 1). It provided a clean layout, a good readability and helps the portfolio looking more calm and professional.} 


\paragraph {I organized my online portfolio into main sections (Introduction, About me, Skills, Projects and Contacts); to make it  easy to navigate and avoid overwhelming the visitors of my website when they are visiting my page, I choose to put a button for my CV, my email account, github account and linkedin account; it helps to find any information my visitors needs quickly without any confusion.}

\paragraph{For the design of my CV, I chose a built-in template in Overleaf and completed it with my informations using La\TeX{}. The choice of the template was based on the German style for CV , and typically the colors used for my website were the same for my CV, to bring in more transparency and consistency. I completed the CV from who I am to my contacts,my languages capacity, professional and soft skills, experience, certification and hobbies to show that I'm not only building myself on hard and professional skills but that outside my coder-bubble I'm a person with interests and dreams.}


\paragraph{As Steve Jobs said \textit{"Design is not just what it looks like and feels like. Design is how it works"}, I tried to maintain transparency, authenticity and professionalism throughout all the designs decisions I've made.}








\newpage

%---------------- Tools and technologies ----------------%
\bfseries\section{Tools and Technologies}
\bfseries\subsection{Github Repository}

Github provides a structured way to store all files related or not to my portfolio. I used it to manage my online portfolio or website files, my Latex CV, my report and some technical exercises. This helped me to ensure transparency and a professional look according to what I was assigned to do.

\bfseries\subsection{My Website}
 HTML(Hyper Text Markup Language) and CSS (Cascading Style Sheet) are the languages I used to layout the content of my webpage, by learning it using youtube videos and websites like WW3 Schools, I applied the  examples and I did my own website, to make it more transparent for me and easy to adjust as needed.

 \bfseries\subsection{La\TeX{}}

 La\TeX{} was used to create my CV and my report. I chose Latex not only because I was assigned to use it but also because it produces clean and professional formatted text.  Latex also ensures that my CV and report would be consistent on different platforms and devices.

  \bfseries\subsection{Overleaf}
  Overleaf is an online Latex editor that in less words simplifies Latex documents. I used it as recommended because it has so many built-in templates that it became easier for me to create my CV, although I tried to do the same thing with my report, I ended up not using a template but created my own report using what I learned as well as  Youtube videos. This made the conception and writing process faster and more organized. 
\newpage

%---------------- Strength and weaknesses ----------------%
\section{Strengths and Weaknesses }
\subsection{Strengths}

I believe that as strength is a quality or feature that makes something effective or useful (Merriam Webster dictionary, nd), the following paragraphs will explain how I defined the strengths of my portfolio:

\paragraph{\textbf{*Clear and simple navigation}: As my portfolio is divided in only 5 parts (Home, about me, skills, projects and contacts), it makes the navigation easier and let the visitor gets the information they want quickly.}

\paragraph{\textbf{*Minimalist design}: By using only 2 colors(Black and Blue), the portfolio maintains a clear, consistent and professional look.}

\paragraph{\textbf{*Updates and maintainance}: As my portfolio is my own work, it is easy for me to update it and maintain it in any way I want, foe example, changing the fonts, updating my skills and projects or add any important information.}

\paragraph{\textbf{*Ready-to-use portfolio}: I believe that this assignemnt was not only to complete a simple portfolio and get the grades, but I kept my heart into my portfolio, that it can be used directly for professional matters and it has a built-in CV.}

\paragraph{As people say "Nothing is perfect", after the consideration of what I call "strengths" of my portfolio, weaknesses also appeared.}

\bfseries\subsection{Weaknesses}

Weaknesses are particular qualities of something that are not good or not effective (Cambridge dictionary, nd), and the following paragraphs explain what I find as weaknesses for my portfolio:

\paragraph{\textbf{*Lack of interactivity}: As I have limited coding experience in HTML/CSS, my portfolio is static and not dynamic, as I couldn't do the right code to add vibrant features and dynamic contents to make it more creative as I am.}

\paragraph{\textbf{*Presented projects}: Due to limited knowledge in Python, the portfolio contains small and beginner python projects, and those can be a constraint if I want to apply and be recruited as software engineer intern.}

\paragraph{\textbf{*Images and screenshots}: I believe that the portfolio could be improved by adding visuals (for projects especially), to make the content more creative and informative.}

\paragraph{As \textit{Marc Jacobs} said:  \textit{ "Always find beauty in things that are odd and imperfect- they are much more interesting."}}




\newpage

\bfseries\section{Recommendations }
A portfolio is your online presence, the guide to whom you are before even talking or even recruiters seeing you, and as important as it is, future recommendations would not miss.

My recommendations for my future portfolio are:

\begin{itemize}
 \item {Add more projects: Developing the portfolio with more professional and advanced projects to showcase my commitment and my skills.}

 \item {Include dynamic features: Adding dynamic features and creative visuals would enhance the look of the portfolio and make it more engaging. }

\item {Improve HTML/CSS coding skills: I believe that if I had enough HTML/CSS knowledge, my portfolio would be different, I would have added more features and accessibility option.}

\item {Add a learning section and recommandation section: I have a vision that my future portfolio will have a learning and recommendation section, so my visitors can see what I'm learning, what I've learned and what I'm intending to learn with also recommandations from peers, colleagues or even chefs to gain more trust and credibility.} 

\end{itemize}

\paragraph{Future features will be added for sure, but I'm grateful for what I have now}

 

 

\bfseries\section{Conclusion}
This portfolio was not only an assignement to complete, but it was a complete learning journey for me. From building from scratch my portfolio and my report, I gained valuable skills not only as an aspiring software engineer but also soft skills such as time management, reliability and self-trust.

\paragraph {This report was built in many stages including a part for the links of my portfolio, technical exercises, Latex CV and report, afterwards the structure of my portfolio and Github repository was particularly explained showing how my portfolio is built and how many repositories I have in my Github, afterwards, I continued by explaining my design decisions, why I chose what I chose and which factors I've considered to built it, and then the tools and technologies I used took place; explaining which technologies I used and tried to built not only my portfolio but also my Latex Cv and my report, and I ended my report by stating what I found as strengths and weaknesses of my portfolio as well as the recommandations I could consider in the near future.}


\paragraph {This journey was not easy, but I'm glad I was able to build and complete it, it has built me and shaped my personality and skills in a way I was not expecting, I was able to grow in knowledge and professionalism and as Logan Aldridge said: \textit{ " We grow through what we go through". }}

\newpage

%---------------- Reference ----------------%
\bfseries\section{Reference}
\paragraph {Bennett, K. (nd). “ #TBT: History of the job search”. \textit{Critical Research}. Available at: \href{https://criticalresearch.com/technology/tbt-history-of-the-job-search/ }(Accessed  22.06.2026)}

\paragraph {Cait, C. (2023). “Electronic portfolio”. \textit{EBSCO}. Available at:  \href {https://www.ebsco.com/research-starters/business-and-management/electronic-portfolio} (Accessed 20.06.2026)}

\paragraph {Cerrato, H. (2012). “The meaning of colors”. Available at:  \href  {https://blocs.xtec.cat/gemmasalvia1617/files/2017/02/the-meaning-of-colors-book.pdf} (Accessed 24.06.2026)}

\paragraph {Learning portal. (2022). "What is a digital portfolio". Available at:\href  {https://tlp-lpa.ca/digital-skills/digital-portfolio} (accessed 24.06.2026)}



\paragraph {Insler, C. (2026). “The 1900 Invention that still controls your career.”.\textit{ Medium}. Available at: \href {https://medium.com/the-managers-mind-a-history-of-work/the-1900-invention-that-still-controls-your-career-0a5444a3ae23} (Accessed 23.06.2026)}

\paragraph {Merriam Webster Dictionary. (nd). “Strength/ Weakness”. Available at: \href {https://www.merriam-webster.com/dictionary/strength} (Accessed 24.06.2026)}



\newpage



\bfseries\section{Appendices}
\subsubsection{Appendix 1}

Vaiana showing the resilience of a young lady protecting her village and exploring a new environment by going beyond her fears and made possible what seemed to be impossible.
 
 \includegraphics[width=0.40\textwidth]{vaiana.jpg}\par\vspace{1cm}


\subsubsection{Appendix 2}
We have a minor inconvenience for the link for the technical exercises, when you run it, it doesn't work but when you add a "-" in "My-python-projects" it runs. I believe it's a compilation problem, but when you are directed to my github you can see all the repositories including this one also. 


\end{document}
